%!TEX root = lebenslauf.tex
\subsection*{Lehre}
\ifshowLehreBetreuung
\begin{quote}
	% Thomas Gorr -15.12.2015
 	\begin{tabular}{p{4cm}l}
 		\multicolumn{2}{l}{Betreuung von Abschlussarbeiten}\\\\
 		2016&Bachelorarbeit\\
 		&\qquad Indoorlokalisierung mit mobilen Ger"aten \\
 		&\qquad und Bluetooth Low Energy Beacons\\\\
 		&\begin{minipage}{10cm}
 		Diese Bachelorarbeit untersucht die Machbarkeit einer Indoorlokalisierung zur ortsbezogenen Kennzahlenanalyse am Beispiel eines Industrieunternehmens. Nach einem Vergleich verf"ugbarer Indoor-Lokalisierungstechnologien im Zusammenspiel zwischen mobilen Ger"aten und Bluetooth Low Energy Beacons wird eine Messdaten-App für mobile Ger"ate entwickelt. Diese dient dem Vergleich verschiedener Bluetooth Low Energy Beacons mit unterschiedlichen mobilen Ger"aten auf variablen Distanzen. Auf Basis der Messdaten-App wird dann eine Kennzahlen-App entwickelt, die prototypisch zur ortsbezogenen Kennzahlenanalyse eingesetzt wird.
		\end{minipage}\\\\\\
	\end{tabular}
	% Nils Brinker -15.12.2015
 	\begin{tabular}{p{4cm}l}
 		2015&Bachelorarbeit\\
 		&\qquad Konzeptuelle Gamifizierung von Software\\
 		&\qquad Entwicklung des CGS-Models\\\\
 		&\begin{minipage}{10cm}
 		Obwohl Gamification in den letzten Jahren einen enormen Aufschwung erlebte, existieren kaum Modelle zur konzeptuellen Entwicklung gamifizierter Software. Deshalb widmet sich diese Arbeit der Entwicklung des CGS-Models (Conceptual Gamification of Software). Es soll Auftraggebern und Entwicklern als Vorgehensweise zur Konzeption gamifizierter Anwendungen dienen, indem es die Entwickler schrittweise und kontextunabh"angig zur Schaffung eines Spielidee- und Spielablauf-fokussierenden Konzepts f"uhrt.
		\end{minipage}\\\\\\
	\end{tabular}
	% Andreas Kaiser - 02.12.2015
 	\begin{tabular}{p{4cm}l}
 		2015&Bachelorarbeit\\
 		&\qquad Integration kontinuierlicher Ortsbestimmung\\
 		&\qquad und Geofencing in Android Apps\\\\
 		&\begin{minipage}{10cm}
 		Diese Arbeit befasst sich mit der Konzeption und Entwicklung einer mobilen Anwendung, die dem Nutzer standortbezogen Inhalte wie z. B. Fotos pr"asentieren kann. GPS-Informationen und Sensordaten des mobilen Endger"ats werden miteinander verkn"upft, um Richtung und Distanz zu verschiedenen interessanten Orten anzuzeigen. Eine Schwierigkeit hierbei ist, kontinuierlich und akkurat die Position des Anwenders zu bestimmen, den Akku jedoch m"oglichst wenig zu belasten.
		\end{minipage}\\\\\\
	\end{tabular}
 % Jan-Nicklas Troschke - 25.09.2015
 	\begin{tabular}{p{4cm}l}
 		&Bachelorarbeit\\
 		&\qquad Integration eines Frameworks f"ur UI-Testing\\
 		&\qquad von Android Anwendungen in ein \\
 		&\qquad  E-Assessment-System\\\\
 		&\begin{minipage}{10cm}
 		Das E-Assessment-System JACK wird an der Universit"at Duisburg-Essen zur Unterst"utzung der Lehre genutzt. Die automatische Auswertung von Programmieraufgaben kann Lehrende entlasten und Studierenden individuelle R"uckmeldung zu ihren L"osungen liefern. "Ubungsaufgaben zur Programmierung f"ur mobile Systeme k"onnen bisher nicht in JACK "uberpr"uft werden. Daher wird in dieser Arbeit ein Testframework f"ur Android Anwendungen ausgew"ahlt und integriert. Auf Basis der Integration wird untersucht, welche Probleme beim Test von mobilen Anwendungen auftreten.
		\end{minipage}\\\\\\
	\end{tabular}
% Ersan Albay - 30.1.2015
	\begin{tabular}{p{4cm}l}
		&Bachelorarbeit\\
		&\qquad Entwicklung eines Konzeptes zum Einsatz\\
		&\qquad von Smartphones in der Lehre\\\\
		&\begin{minipage}{10cm}
		Die Bachelorarbeit ist in den Kontext technologische Unterst"utzung von Classroom Response Systemen einzuordnen. Der Autor setzt sich mit der Problematik auseinander, dass Vorlesungsveranstaltungen an Universit"aten mit vielen Teilnehmern zu einer interaktionslosen, einseitig gepr"agten Kommunikation zwischen Dozenten und Studierenden degenerieren. W"unschenswert w"are hingegen eine (inter-)aktive Einbeziehung der Studierenden. Der Autor untersucht in seiner Bachelorarbeit wie Smarthphones und Tablet Computer verwendet werden k"onnen, um Vorlesungsveranstaltungen interaktiver zu gestalten.
		\end{minipage}\\\\\\
	\end{tabular}
%
% Chinh Truc Dao - 26.1.2015
	\begin{tabular}{p{4cm}l}
		&Masterarbeit\\
		&\qquad Anpassung des Android OS zur\\
		&\qquad Erm"oglichung kontextsensitiver Anwendungstests\\
		&\qquad au"serhalb emulierter Umgebungen\\\\
		&\begin{minipage}{10cm}
		Die Arbeit ist den Kontext Testen mobiler Anwendungen einzuordnen. Das Testen mobiler Anwendungen unterliegt in besonderem Ma"se den durch die technischen Grundlagen der jeweiligen mobilen Plattform vorgegebenen Einschr"ankungen. Gegenw"artige Werkzeugunterst"utzung mobiler Plattformen f"ur die Anwendungsentwicklung und insbesondere f"ur den Anwendungstest ist nur rudiment"ar ausgepr"agt, so dass wesentliche f"ur mobile Plattformen charakteristische Funktionen lediglich manuell getestet werden k"onnen. Das trifft insbesondere zu auf sensorbasierte Anwendungen, die Sensordaten als zus"atzlichen Eingabekanal verwenden. Neben dem Fehlen ad"aquater Schnittstellen und automatisierter Testwerkzeuge stellt sich hier bei der manuellen Testdurchf"uhrung die Herausforderung, dass menschliche Tester nicht in der Lage sind physikalische Bedingungen der Testumgebung zu "uberpr"ufen und gezielt bereitzustellen. Gegenstand der  Masterarbeit ist deshalb die Konzeption und Implementierung einer technischen Schnittstelle im Android-Betriebssystem, durch welche Testen kontextsensitiver Anwendungen (d.h. von Anwendungen, u.a. physikalische Parameter der Betriebsumgebung als Eingabekanal verwenden) auf mobilen Ger"aten zur Laufzeit durch Simulation von Sensordaten erm"oglicht wird. Die Komponenten des Betriebssystems wurde so modifiziert, dass zur Laufzeit extern Sensormesswerte bereitgestellt werden k"onnen. Die implementierte Schnittstelle wurde schlie"slich auf dem Google Nexus S und dem Android-Ger"ateemulator evaluiert.
		\end{minipage}\\\\\\
	\end{tabular}
%
% Marlen Skarabinia - 14.1.2015
	\begin{tabular}{p{4cm}l}
		&Bachelorarbeit\\
		&\qquad Datenschutz im Speditionswesen\\\\
		&\begin{minipage}{10cm}
		Die Arbeit ist in den Bereich mobile Gesch"aftsprozesse an der Schnittstelle zu Unternehmen-Informationssystemen einzuordnen. Sie befasst sich mit dem Thema Datenschutz im Speditionswesen im besonderen Hinblick auf eine Konzepterstellung f"ur ein Flottenmanagementsystem unter Ber"ucksichtigung rechtlicher Rahmenbedingungen zum Datenschutz. Flottenmanagementsysteme bieten neben den rein inhaltlichen Funktionen des Managements von Fuhrparks und der Zuordnung von Ressourcen ebenfalls die theoretische M"oglichkeit der unzul"assigen "uberwachung von Arbeitnehmern. Die Autorin der Bachelorarbeit diskutiert im ersten Teil der Arbeit die aktuelle Lage zum Besch"aftigtendatenschutz. Im zweiten Teil der Arbeit werden ausgew"ahlte Flottenmanagementsysteme anhand ihrer Systembeschreibungen auf Konformit"at gepr"uft. Im Anschluss erstellt die Autorin ein Konzept f"ur ein verbessertes neues Flottenmanagementsystem.
		\end{minipage}\\\\\\
	\end{tabular}
%
% Ey"up Poyraz  - 2.12.2014
	\begin{tabular}{p{4cm}l}
		2014&Bachelorarbeit\\
		&\qquad Bewertung alternativer Implementationstechnologien\\
		&\qquad  f"ur eine mobile App im Facility-Management\\\\
		&\begin{minipage}{10cm}
		Die Arbeit ist den Bereich Mobilisierung von Gesch"aftsprozessen durch Verwendung mobiler Anwendungen (sogenannter Apps) auf mobilen Ger"aten (Smartphone und Tablet Computer) im Wirtschaftszweig Facility-Management einzuordnen. Basierend auf dem expandierenden Markt f"ur mobile Ger"ate und gleichzeitiger Divergenz der jeweiligen Technologien setzt sich der Autor mit der Frage auseinander, ob native Technologien sogenannten Cross-Platform-Technologien bei der Implementierung typischer Aufgaben im Facility-Management "uberlegen sind oder ob Cross-Platform-Technologien technologisch gleichwertig native L"osungen ersetzen k"onnen. Gegenstand der Untersuchung die ist Anwendung Handwerkerkopplung. Diese erm"oglicht die papierlose Abwicklung von Reparaturauftr"agen, Abrechnung und Dokumentation im Facility-Management Gegenw"artig wird die Anwendung von der Firma BAGUS GmbH als Web-Anwendung implementiert. Diese Bachelorarbeit untersucht wissenschaftliche, technische und wirtschaftliche Aspekte der Erg"anzung der Web-Anwendung durch eine mobile Anwendung auf der Basis von Cross-Platform-Technologien.
		
		\end{minipage}\\\\\\
	\end{tabular}
%
% Ilja Averdieck  - 25.11.2014
	\begin{tabular}{p{4cm}l}
		&Bachelorarbeit\\
		&\qquad Erarbeitung eines interaktiven Lehrkonzepts\\
		&\qquad  unter Einsatz von Wearables\\\\
		&\begin{minipage}{10cm}
		Die Bachelorarbeit ist in den Kontext technologische Unterst"utzung von Classroom Response Systemen einzuordnen. Der Autor erarbeitet ein Konzept f"ur eine Unterst"utzung des Lehrenden durch mobile Ger"ate, insbesondere Wearables, in einen interaktiven Lernumfeld. Frontale Lehrmethoden schaffen ein Unterrichtsklima, in dem es f"ur individuelle Studierende kaum einen Anreiz gibt, sich bereits w"ahrend des Unterrichts interaktiv mit dem Thema auseinanderzusetzen. Insbesondere existiert wenig Anreiz auf Fragen des Lehrenden oder Dozenten zum Verst"andnis des Lehrstoffs "offentlich zu antworten. Um einen solchen Anreiz zu schaffen, wurde eine mobile Variante eines Classroom Response Systems konzipiert und prototypisch umgesetzt, die es dem Lehrenden gestattet interaktive Element in den Unterricht zu integrieren und Studierenden die M"oglichkeit einzur"aumen anonym zu interagieren. Um gleichzeitig zu gew"ahrleisten, dass der Lehrende nicht durch die Bedienung technischen Ger"ats abgelenkt wird, kommen sogenannte Wearbables in Form von Smartwatches zur Anwendung. 
		
		\end{minipage}\\\\\\
	\end{tabular}
%
% Alexander Massolle - 25.11.2014
	\begin{tabular}{p{4cm}l}
		&Masterarbeit\\
		&\qquad Informationsgewinnung aus nat"urlicher Sprache\\
		&\qquad zur Interaktion mit mabilen Anwendungen\\\\
		&\begin{minipage}{10cm}
		Die Arbeit ist in den Bereich Software Engineering f"ur mobile Systeme an der Schnittstelle zu Human Computer Interaction einzuordnen. Der Autor betrachtet eine L"osung der Fragestellung, wie die nat"urliche Sprache als Eingabevektor zur Steuerung einer mobiler Anwendung verwendet werden kann. Der Autor erkennt Bedarf an der M"oglichkeit mobile Anwendungen mit nat"urlicher Sprache bedienen zu k"onnen, um Anwendungen auch in solchen Umgebungen nutzbar zu machen, in denen auf prim"are Eingabemethoden mobiler Ger"ate -- n"amlich Bedienung durch Ber"uhrung des Displays mit einem oder mehreren Fingern -- verzichtet werden muss, beispielsweise beim F"uhren einen Kraftfahrzeugs. Betrachtet werden die Grundlagen der Signalverarbeitung zur Informationsgewinnung aus mit einem Mikrofon aufgezeichneter nat"urlicher Sprache, die Transformation zu Text in einem maschinenlesbaren Format, die Extraktion verwertbarer Information aus der textuellen Repr"asentierung und die Integration dieser Information in die Interaktion eines Anwenders mit einer Software. 
		\end{minipage}\\\\\\
	\end{tabular}
%
% Marc Gesth"usen - 21.7.2014
	\begin{tabular}{p{4cm}l}
		&Bachelorarbeit\\
		&\qquad Testautomatisierung sensorbasierter mobiler\\
		&\qquad Anwendungen mit Calabash-Android\\\\
		&\begin{minipage}{10cm}
		Die Arbeit ist in den Bereich Software Engineering f"ur mobile Systeme einzuordnen. Sie befasst sich mit Automatisierungstechnologien zur Unterst"utzung der Testdurchf"uhrung f"ur sensorbasierte mobile Anwendungen f"ur die Plattform Android. Betrachtet wird prim"ar das Testautomatisierungsrahmenwerk Calabash f"ur die Plattform Android. Dieses wird durch Integration einer erweiterten Variante des Werkzeugs OpenIntents SensorSimulator so aufgewertet, dass eine Realisierung von Tests f"ur solche Anwendungen erm"oglicht wird, die Sensordaten als Eingabevektoren verwenden.
		\end{minipage}\\\\\\
	\end{tabular}
%
% Rinay Baron - 17.7.2014
	\begin{tabular}{p{4cm}l}
		2014&Bachelorarbeit\\
		&\qquad Frontend Performance Test\\
		&\qquad eines Smart Home Systems\\\\
		&\begin{minipage}{10cm}
		Die Arbeit ist in den Bereich Testen dynamischer Web-Anwendungen einzuordnen. Die Autorin Untersucht die Fragestellung, wie das Web-basierte Frontend-System einer Smart Home L"osung automatisiert und mit geringem personellen Aufwand hinsichtlich nichtfunktionaler Anforderungen, insbesondere Performanz, getestet werden kann. Motiviert wird die Problemstellung  basierend auf der Erkenntnis, dass ein wesentlicher Teil der Gesamtladezeit einer Web-Anwendung auf das Frontend entf"allt und Verz"ogerungen deshalb als besonders st"orend vom Anwender wahrgenommen werden. Der in dieser Bachelorarbeit entwickelte Frontend Performance Test soll die Grundlage zur Optimierung der Gesamtladezeit der betrachteten Smart Home L"osung sein.
		
		\end{minipage}\\\\\\
	\end{tabular}
%
% Michael Roedel - 22.5.2014
	\begin{tabular}{p{4cm}l}
		&Bachelorarbeit\\
		&\qquad Analyse von Gesch"aftsprozessen hinsichtlich\\
		&\qquad ihrer Eignung zur zur Unterst"utzung durch\\
		&\qquad mobile Systeme mit Hilfe eines\\
		&\qquad Kriterienkatalogs\\\\
		&\begin{minipage}{10cm}
		Die Arbeit ist in den Bereich Prozessanalyse zur Mobilisierung von Gesch"aftsprozessen einzuordnen. Der Autor behandelt die Frage, welche Gesch"aftsprozesse sich unter welchen Bedingungen zur Unterst"utzung durch mobile Software eignen. Hierzu werden Kriterien erstellt, anhand derer Gesch"aftsprozesse basierend auf Aufwand und Nutzen hinsichtlich ihres Mobilisierungspotenzials bewertet werden k"onnen. Eine Vorgehensweise erm"oglicht basierend auf diesen Kriterien das Auffinden solcher Gesch"aftsprozesse in digitalen Unternehmen, die sinnvoll durch mobile Software unterst"utzt werden k"onnen.
		\end{minipage}\\\\\\
	\end{tabular}
%
% Martin Hermes - 14.1.2014
	\begin{tabular}{p{4cm}l}
		&Masterarbeit\\
		&\qquad Positionsbasierte Paarerkennung\\
		&\qquad zwischen bewegten Objekten\\\\
		&\begin{minipage}{10cm}
		Die Arbeit untersucht das Thema der positionsbasierten Paarfindung. Insbesondere wird der Sonderfall der Paarerkennung zwischen Personen und Fahrzeugen betrachtet, der trotz vielseitiger Einsatzm"oglichkeiten bisher nicht tiefergehend untersucht wurde. Es werden L"osungen vorgestellt, die es erlauben die Paarerkennung zuverl"assig, performant und effizient durchzuf"uhren.
		\end{minipage}\\\\\\
	\end{tabular}
%
% Don Qui Ho - 16.3.2013
	\begin{tabular}{p{4cm}l}
		2013&Masterarbeit\\
		&\qquad Entwicklung eines Toolkits zur Unterst"utzung\\
		&\qquad von automatisiertem Testen von Android Apps\\
		&\qquad mit sensorenbasierter Gestenerkennung\\\\
		&\begin{minipage}{10cm}
		Die Arbeit untersucht eine Erweiterung des Open-Source Projektes \emph{SensorSimulator} um die M"oglichkeit, mit Hilfe von JUnit automatisierte Tests von Sensoren-unterst"utzten Android Apps durchzuf"uhren. Der bestehende Code wurde um Schnittstellen erweitert, welche eine Einbindung in JUnit Testprojekte erm"oglichen. Beliebige Sensorwerte und Sequenzen davon sollen als Testparameter eingestellt werden k"onnen, um die Reaktion der App darauf zu testen und mit "ublichen Methoden zu verifizieren.
		\end{minipage}\\\\\\
	\end{tabular}
%
% Frank Schnellenkamp - 17.8.2012
	\begin{tabular}{p{4cm}l}
		2012&Bachelorarbeit\\
		&\qquad Testautomatisierung von Android Anwendungen\\\\
		&\begin{minipage}{10cm}
Die Bachelorarbeit untersucht und bewertet Testautomatisierungs-Frameworks f"ur die mobile Plattform Android hinsichtlich ihrer Eignung zur Erstellung und Durchf"uhrung vollst"andig automatisierter funktionaler Anwendungstests. Ausgehend von einer Beispielanwendungen werden Testf"alle und Untersuchungskriterien Entworfen anhand derer eine Diskussion existierender Testautomatisierungs-Frameworks erfolgt.
		\end{minipage}\\\\\\
	\end{tabular}
%
%
\end{quote}
\fi
\ifshowLehreEssen
\begin{quote}
	\begin{tabular}{p{4cm}l}
		\multicolumn{2}{l}{Universit"at Duisburg-Essen}\\\\
		\multicolumn{2}{l}{Betreuung Vorlesungen}\\\\
		2012 - 2016&Software Engineering f"ur mobile Systeme\\\\
		&\begin{minipage}{10cm}
			Inhaltliche und organisatorische Erarbeitung sowie Durchf"uhrung der Vorlesung Software Engineering f"ur mobile Systeme. Konzeption, Planung und Durchf"uhrung einer begleitenden Seminarveranstaltung.\\\\
			Die Vorlesung vermittelte einen allgemeinen "Uberblick "uber das Thema Software-Engineering f"ur mobile Systeme und betrachtete dabei insbesondere die Grundlagen der Entwicklung mobiler Systeme, Grundlagen der Programmierung mobiler Benutzerschnittstellen, Grundlagen plattformspezifischer Konzepte, Backend-basierte Anwendungen / Kommunikation mit Services, plattformspezifische Vertiefung der Programmierung und Cross-Plattform Strategien und Technologien. Neben diesen grundlegenden Konzepten und Prinzipien, wurden  ausgew"ahlte mobile Plattformen vorgestellt. Praktische Kenntnisse wurden intensiv am Beispiel der Plattform Android erarbeitet.
		\end{minipage}\\\\\\
		2013 - 2015&Modelle der Informatik\\\\
		&\begin{minipage}{10cm}
		Inhaltliche und administrative Betreuung der Vorlesung Modelle der Informatik. Regelm"a"sig durchf"uhrung der Vorlesung. Vermittlung von Modellierungsparadigmen und Formalismen, die sich in der praktischen Anwendung als erfolgreich erwiesen haben; Erlernen der grundlegenden Konzepte, der zugeh"origen Formalismen und Notationen, der Anwendungsbereiche und der wichtigsten Algorithmen der Themengebiete Formale Sprachen, Logik, B"aume, Graphen und Netzwerke, Petri-Netze und Objektorientierte Modellierung mit Unified Modeling Language (UML).
		\end{minipage}\\\\\\
	\end{tabular}
	\begin{tabular}{p{4cm}l}
		2011&Spezifikationssprachen\\\\
		&\begin{minipage}{10cm}
			Betreuung der Vorlesung Spezifikationssprachen und Durchf"uhrung von Seminarveranstaltungen.\\\\
			Gegenstand der Vorlesung waren sowohl ein allgemeiner "Uberblick "uber das Thema Spezifikation von Software-Systemen als auch die Betrachtung verschiedener Spezifikationssprachen. Neben grundlegenden Konzepten und Prinzipien von Spezifikationssprachen wurden ausgew"ahlte Sprachen detailliert vorgestellt.
		\end{minipage}\\\\\\
		\multicolumn{2}{l}{Betreuung Seminare}\\\\
		2011-2015&Projektseminar Software Engineering\\\\
		&\begin{minipage}{10cm}
		Thematische und administrative Betreuung von Studenten in Seminarveranstaltungen zu Themen des angewandten Software Engineering.
		\end{minipage}\\\\\\
	\end{tabular}
	\begin{tabular}{p{4cm}l}
		\multicolumn{2}{l}{Universit"at Leipzig}\\\\
		\multicolumn{2}{l}{Betreuung Vorlesungen}\\\\
		2009-2010&Ausgew"ahlte Kapitel der Softwaretechnik\\\\
		&\begin{minipage}{10cm}
			Unterst"utzung der Vorlesung, Durchf"uhrung einer Seminarveranstaltung
		\end{minipage}\\\\\\
	\end{tabular}
\end{quote}
\fi